\section{Evolusi Mikroprosesor dan Kedudukan 8086}
\label{sec:evolusi-8086}

Mikroprosesor adalah satu chip silikon yang memuat seluruh unit
pengolah utama komputer. Seluruh program assembly yang ditulis pada
mata kuliah ini pada akhirnya diterjemahkan menjadi deretan sinyal
yang mengalir di dalam chip tersebut. Karena itu, memahami arsitektur
mikroprosesor berarti memahami ``peta jalan'' yang dilalui setiap
instruksi, sejak diambil dari memori sampai hasilnya ditulis kembali.

\subsection{Dari 4004 ke 8086}

Sejarah mikroprosesor dimulai pada awal dekade 1970-an. Setiap
generasi menambah lebar kata (\eng{word}), lebar bus alamat, dan
jumlah transistor, sehingga kemampuan pengalamatan memori meningkat
pesat. Tabel~\ref{tab:evolusi} merangkum tonggak penting yang
relevan bagi mata kuliah ini.

\begin{table}[H]
    \centering
    \caption{Tonggak evolusi mikroprosesor menuju 8086.}
    \label{tab:evolusi}
    \small
    \begin{tabular}{@{}clccp{4.0cm}@{}}
        \toprule
        \textbf{Tahun} & \textbf{Prosesor} & \textbf{Bus data} &
        \textbf{Bus alamat} & \textbf{Catatan} \\
        \midrule
        1971 & Intel 4004 & 4 bit & 12 bit &
        Mikroprosesor komersial pertama, untuk kalkulator \\
        \addlinespace
        1974 & Intel 8080 & 8 bit & 16 bit &
        Ruang alamat 64 KB; dasar microcomputer awal \\
        \addlinespace
        1978 & Intel 8086 & 16 bit & 20 bit &
        Ruang alamat 1 MB; dasar keluarga x86 \\
        \addlinespace
        1979 & Intel 8088 & 8 bit (luar) & 20 bit &
        Inti 16-bit, bus luar 8-bit; dipakai IBM PC \\
        \addlinespace
        1982 & Intel 80286 & 16 bit & 24 bit &
        Menambah \eng{protected mode} dan 16 MB alamat \\
        \bottomrule
    \end{tabular}
\end{table}

Dua hal penting terlihat dari Tabel~\ref{tab:evolusi}. Pertama,
8086 muncul pada 1978 dengan bus data 16-bit dan bus alamat 20-bit.
Kedua, varian 8088 yang dipakai IBM PC mempertahankan inti 16-bit
yang sama, hanya menyempitkan bus luar menjadi 8-bit agar cocok
dengan rangkaian pendukung yang sudah ada saat itu.

\subsection{Mengapa 8086 Masih Dipelajari}

Meskipun usianya sudah lebih dari empat dekade, 8086 tetap dipakai
sebagai model pembelajaran karena tiga alasan
\cite{intel1979,stallings2016}.

\begin{enumerate}
    \item \textbf{Sederhana namun lengkap.} 8086 sudah memiliki
          seluruh konsep penting prosesor modern: register, bus,
          \eng{flag}, segmentasi, dan \eng{pipeline} sederhana,
          tanpa mekanisme yang rumit seperti \eng{cache} berlapis.

    \item \textbf{Mudah diamati.} Emulator seperti Emu8086
          memperlihatkan isi register, memori, dan \eng{flag} secara
          langsung pada setiap langkah eksekusi, sehingga konsep
          abstrak dapat diverifikasi sendiri oleh mahasiswa.

    \item \textbf{Akar keluarga x86.} Seluruh prosesor x86 modern
          tetap kompatibel ke belakang. Sebuah program 8086 mode real
          masih dapat dijalankan pada prosesor masa kini melalui mode
          kompatibilitas \cite{carleton2011}.
\end{enumerate}

\subsection{Arsitektur Internal dan Arsitektur Eksternal}

Pembahasan arsitektur mikroprosesor dibedakan menjadi dua sudut
pandang yang harus dipisahkan dengan jelas.

\begin{definisi}[Arsitektur internal]
\term{Arsitektur internal} (\eng{internal architecture}) adalah
susunan unit fungsional di dalam chip, yaitu register, \eng{ALU},
unit pengendali, dan saluran data internal. Arsitektur ini menjelaskan
\emph{bagaimana} prosesor mengolah sebuah instruksi.
\end{definisi}

\begin{definisi}[Arsitektur eksternal]
\term{Arsitektur eksternal} (\eng{external architecture}) adalah
tampak luar chip, yaitu jumlah dan fungsi pin, kelompok bus, serta
sinyal kendali yang menghubungkannya dengan memori dan perangkat
I/O. Arsitektur ini menjelaskan \emph{bagaimana} prosesor berbicara
dengan dunia luar.
\end{definisi}

Kedua sudut pandang tersebut saling melengkapi. Seorang programmer
assembly lebih sering berhadapan dengan arsitektur internal karena
ia menuliskan nama register. Sebaliknya, seorang perancang sistem
digital lebih banyak bekerja pada arsitektur eksternal karena ia
harus menyambungkan pin \reg{A0}--\reg{A19}, \reg{AD0}--\reg{AD15},
dan sinyal kendali ke rangkaian lain.

\begin{infobox}
Kesalahan pemahaman yang sering muncul adalah menganggap register
sebagai ``kotak di memori''. Register sebenarnya berada di dalam
chip, bukan di memori, dan jumlahnya sangat sedikit. Bab ini membahas
kedua sisi arsitektur tersebut secara berurutan: sisi internal pada
Bagian~\ref{sec:biu-eu}, sisi eksternal pada Bagian~\ref{sec:bus-siklus}.
\end{infobox}

\subsection{Kotak Hitam: Cara Pandang Sederhana}

Sebelum masuk ke diagram blok yang rinci, prosesor dapat dipandang
sebagai sebuah kotak hitam. Ke luar, kotak itu hanya berhubungan
melalui empat kelompok saluran: bus alamat 20-bit, bus data 16-bit,
bus kontrol, dan sumber catuan serta pewaktu. Setiap instruksi yang
dijalankan pasti melibatkan pertukaran informasi melalui saluran
tersebut. Isi kotak itulah yang dibahas pada bagian berikutnya.
